%\usepackage [ISO-8859-1]{inputenc}
\documentclass[a4paper,12pt,oneside]{amsbook}
\usepackage[T1]{fontenc}
\usepackage{listings}
\usepackage{textcomp}
\usepackage{multicol,calc}
%\usepackage{euler}
%\usepackage{garamond}
\usepackage{hyperref}

\usepackage{newcent}
%\fontfamily{fxm}\selectfont
\usepackage{fancyhdr}
\pagestyle{fancy}
% Detta paket skoter grafiken
%\usepackage[dvips]{graphicx}
\usepackage[pdftex]{graphicx}

\usepackage{amsthm}
\usepackage{amsmath}
\usepackage[swedish]{babel}
% For bortkommentering
\usepackage{verbatim}
% For att slippa indentering av Sections och SubSections
\usepackage{indentfirst}

% Har stalls paragrafindentering och avstand mellan stycken och rader mm
\setlength{\parindent}{0pt}
\setlength{\parskip}{6pt}
\setlength{\baselineskip}{12pt}
% Har definieras sidan
\setlength{\voffset}{-15mm}
\setlength{\hoffset}{-5.5mm}

\setlength{\topmargin}{0mm}
\setlength{\headheight}{10mm}
\setlength{\headsep}{8mm}
\setlength{\footskip}{16mm}

\setlength{\evensidemargin}{0mm}
\setlength{\oddsidemargin}{0mm}

\setlength{\marginparsep}{0mm}
\setlength{\marginparwidth}{0mm}

\setlength{\textwidth}{170mm}
\setlength{\textheight}{230mm}
\setlength{\paperwidth}{210mm}
\setlength{\paperheight}{297mm}
\setlength{\headwidth}{170mm}
% Huvud och Fot
%\fancypagestyle{}{
\fancyhf{}\fancyhead[c]{\small{\textit{Programmeringsolympiaden Kvalificering 2024}}}
\fancyfoot[c]{\thepage}
\renewcommand{\headrulewidth}{0.4pt}
\renewcommand{\footrulewidth}{0.4pt}
\newenvironment{lista}
{\begin{itemize}
\setlength{\parindent}{0pt}
\setlength{\itemsep}{6pt}}
{\end{itemize}}
% Uppgifter
\newcounter{probnr}
\newenvironment{tal}{%
\begin{list}
%{\textbf{\arabic{section}.\arabic{probnr}}} {\usecounter{probnr}
{\textbf{\arabic{probnr}}} {\usecounter{probnr}
\setlength{\leftmargin}{0mm}
\setlength{\rightmargin}{0mm}
\setlength{\labelwidth}{-1mm}
\setlength{\labelsep}{1mm}}
\setlength{\itemsep}{6pt}
}{\end{list}}
% Definition av avsnitt: FAKTA, EXEMPEL, PASTAENDE
\newtheorem{fakta}{Fakta}
\newtheorem{exempel}{Exempel}
\newtheorem{problem}{Problem}
%
\newtheoremstyle{test}% NAME
{20pt}      % ABOVESPACE
{10pt}      % BELOWSPACE
{\sffamily} % BODYFONT
{0pt}       % INDENT
{\scshape}  % HEADFONT
{}          % HEADPUNCT
{\newline}  % HEADSPACE
{}          % CUSTOM-HEAD-SPEC

\theoremstyle{test}
\newtheorem{program}{Program}
\newcommand{\sv}[1]{\textsc{#1}}            % Sma versaler
\newcommand{\fe}[1]{\textbf{#1}}            % FET
\newcommand{\ku}[1]{\textit{#1}}            % KURSIV
\newcommand{\cu}[1]{\texttt{#1}}            % Courier
\newcommand{\sk}[1]{\texttt{#1}}            % Courier
\newcommand{\rubrik}[1]{\begin{center}\sf\huge{#1}\normalsize\rm\end{center}}
\begin{document}
%\DeclareGraphicsExtensions{.jpg,.pdf,.mps,.png,.eps}



\thispagestyle{fancy}
\lstset{basicstyle=\ttfamily,
  breaklines=true}
% ----------------------------------------------------

\begin{center}
\Huge{Programmeringsolympiaden 2023}
\end{center}
\vspace{-1.2cm} 
\specialsection*{Tävlingsregler för skolkvalet}
\begin{lista}
\item Tävlingen äger rum på av skolan bestämt datum under {\bf fyra
    timmar. Ingen förlängning ges för lunch eller raster.} Eleven ska i förväg komma överens med läraren om att använda egen dator eller en som skolan tillhandahåller. I vilket fall som helst måste eleven befinna sig i avtalad lokal på skolan. 
    %{\bf På grund av pandemin kan i undantagsfall tävling ske från annan plats godkänd av läraren, givet tillräcklig övervakning.}
\item Tävlingen består av fem uppgifter som vardera ska lösas genom ett
  datorprogram i valfritt programmeringsspråk.
\item {\bf Indata kan läsas in i programmet på valfritt sätt}, t.ex. genom att programmet för en dialog med användaren (som i körningsexemplena i uppgifterna), att de skrivs in i ett grafiskt gränssnitt eller att datafiler skickas till {\em standard input}. Kom bara överens med din lärare om hur programmet ska testas.
\item Dina lösningar kommer att testköras med förpreparerade
  indata. Varje uppgift testas normalt med 5 testfall, som vardera ger
  1 poäng om ditt program skriver ut korrekt svar inom en exekveringstidsgräns av
  {\bf 3 sekunder}. Ingen test av indata behöver göras, den följer specifikationerna 
i uppgiften.
\item Det är ofta olika begränsningar på de olika testfallen, t.ex.\ storleken på
indata eller andra inskränkningar. Detta anges i uppgiften. {\bf Observera att det kan vara
helt olika svårighetsgrad på en uppgift beroende på dessa skillnader. Det kan därför vara
lättare att få delpoäng på en uppgift som verkar svår än att få full
poäng på en uppgift som verkar lättare.} Informationen om delpoäng är
därför extremt viktig för att planera sin
tävling.
\item Rättningen utförs på samma eller likvärdig dator. Ändringar
  i källkoden tillåts ej efter tävlingen. Om programmet
  inte kan kompileras ges 0 p.\ på uppgiften. 
\item Om något av följande inträffar ger
  det {\em testfallet} 0 poäng, men programmet fortsätter testas med övriga testfall.
\begin{lista}
\item Exekveringstiden överstiger 3 sekunder \vspace{-0.2cm}
\item Exekveringsfel (run time error) \vspace{-0.2cm}
\item Fel svar \vspace{-0.2cm}
\end{lista}

\item Deltagandet är individuellt vilket bland annat innebär att inget utbyte av idéer eller 
filer får ske under tävlingen.  
\item Hjälpmedel: Valfritt skriftligt material, material som finns installerat på datorn samt material som finns tillgängligt på internet. Det är \ku{inte} tillåtet att aktivt kommunicera på internet (t.ex. chatta eller ställa frågor till ett forum) utan endast att söka efter information. Räknedosa är tillåten.
\item Tävlingsbidraget ska lämnas in i form av källkodsfiler (uppg1...uppg1 med passande filtillägg) som läggs på överenskommen plats. Ingen hänsyn tas till andra filer. Var noga 
med att lämna in den korrekta versionen av ditt program.  
\end{lista}

%De högst placerade i kvalet går vidare till finalen där landslagsplatser till BOI i Sverige och IOI i Japan står på spel.
\begin{center}
\Large Lycka till!
\end{center}
%\end{document}

\newpage
\specialsection*{Uppgift 1 -- Reduplikation}

Reduplikation är när man upprepar ett ord flera gånger, ofta för att förstärka eller ändra betydelsen.
I vissa språk som indonesiska är reduplikation väldigt vanligt, och där har det många olika
grammatiska funktioner. Men även i svenska används reduplikation ibland, till exempel är det vanligt
att man säger \texttt{"hej hej"} istället för \texttt{"hej"}.

Rama är en flitig användare av reduplikation när han chattar med sina kompisar. Men det tar lång tid
att skriva samma ord upp till nio gånger. Därför brukar han skriva ordet en gång, åtföljt av en
siffra som visar hur många gånger ordet ska upprepas.

Skriv ett program som läser in ett ord som Rama tänker upprepa (som består av som mest 9 bokstäver)
och hur många gånger det ska upprepas (som mest 9). 
Programmet ska sedan skriva ut ordet fast upprepat rätt antal gånger.


\fe{Poängsättning:} För testfall värda $1$ poäng (utav $5$ möjliga) gäller att antalet gånger ordet ska upprepas är $1$.

\vspace{1cm}
\fe{Körningsexempel 1}
\begin{verbatim}
Ordet ? hej
Antal upprepningar ? 3 

Svar: hejhejhej
\end{verbatim}

\vspace{1cm}
\fe{Körningsexempel 2}
\begin{verbatim}
Ordet ? ha
Antal upprepningar ? 4 
  
Svar: hahahahaha
\end{verbatim}



\newpage
\specialsection*{Uppgift 2 -- Tre i rad}


%\begin{center}
%  \includegraphics[width=12cm]{sifferkrypto.pdf} \\ ~\\
%  \ku{Exempel på missuppfattning som kan ske med Alice nya chiffer.}
%\end{center}

\vspace{1cm}

Matematikern Lenore Oljer har gjort så mycket matte att hon har tröttnat på primtal. 
Hon har nu börjat studera så kallade trevliga tripletter. 
Det är väldigt lätt att skapa en trevlig triplett. 
Börja med ett heltal $a$ som är
större än noll och skapa sedan tripletten $(a,a+1,a+2)$. 
Alla tripletter som skapats på detta sätt är trevliga. 
Några exempel på trevliga tripletter är $(4,5,6)$ och $(15,16,17)$.

För att lära sig mer om trevliga tripletter vill Lenore hitta hur 
många tal som kan skrivas som produkten av
alla talen i en trevlig triplett. Hon kallar dessa tal "treiga tal". 
Några exempel på treiga tal är $24=2 \cdot 3 \cdot 4$
och $336 = 6 \cdot 7 \cdot 8$. Lenore vill nu veta hur många treiga tal det 
finns som är mindre än talet $N$.

Skriv ett program som läser in talet $N$ och skriver ut hur många tal mindre än
$N$ är treiga.

\fe{Poängsättning:}\\
De fem testfallen som vardera ger $1$ poäng har följande begränsningar:
\begin{enumerate}
\item  $\;\;\;\;N \le 25 $    
\item  $\;\;\;\;N \le 80 $ 
\item  $\;\;\;\;N \le 1000 $   
\item  $\;\;\;\;N \le 10^5 $   
\item  $\;\;\;\;N \le 10^9 $   
\end{enumerate}


\vspace{1cm}

\fe{Körningsexempel 1}
\begin{verbatim}
N ? 75

Svar: 3
\end{verbatim}

\fe{Förklaring:} De treiga talen är 6, 24 och 60.

\vspace{1cm}

\fe{Körningsexempel 2}
\begin{verbatim}
N ? 24

Svar: 1
\end{verbatim}

\fe{Förklaring:} Eftersom det treiga talet $24$ inte är mindre än $N$ räknar vi bara 6.

\vspace{3cm}

\fe{Körningsexempel 3}
\begin{verbatim}
N ? 1234

Svar: 9
\end{verbatim}



\newpage
\specialsection*{Uppgift 3 -- Inomhusorientering}

En snabbt växande sport är \ku{inomhusorientering}, där de största tävlingarna lockar tusentals deltagare som springer runt på t.ex.\ ett Campus. Genom arrangörens noga uttänkta avspärrningar blir det en labyrintliknande utmaning att ta sig mellan kontrollerna. Ofta gäller det att optimera användandet av trappor, och det är det vi ska titta på i den här uppgiften. Givet en karta som indelats i $N$ \ku{områden} ska du beräkna det minimala antalet trapplöpningar som orienteraren behöver göra för att besöka de $N$ områdena i tur och ordning (för att ta en kontroll i varje).

Byggnaden innehåller ett antal lodräta trapphus betecknade med bokstäver. Varje område ligger på en viss våning och är kopplat till ett eller flera av trapphusen. Ett trapphus behöver inte vara kopplat till något område på varje våning. Exempelvis kan trappa A gå mellan våning 1 och våning 3 utan att vara kopplad till något område på våning 2, som i figuren nedan. Att förflytta sig en våning upp eller ner i en trappa tar en tidsenhet. Däremot bryr vi oss inte om tiden det tar att röra sig inom ett område.

Programmet ska först läsa in antalet områden $N$, där $3\le N \le 15$. För varje område ska programmet läsa in den våning som området befinner sig på (ett heltal mellan 1 och 100) och en sträng som innehåller ett eller flera tecken, valda bland versalerna A--Z, och som anger vilket eller vilka trapphus området är kopplat till. Programmet ska skriva ut ett tal: den minimala tiden det tar att besöka de $N$ områdena i precis den ordning de angavs i indata, med starten i område $1$ och målet i område $N$. 

\fe{Poängsättning:} I testfall värda $1$ poäng (av $5$ möjliga) kommer alla strängar vara \cu{A}.\\

\begin{multicols*}{3}

\fe{Körningsexempel 2}
\begin{verbatim}
N ? 8
Omr 1, vån ? 1
Omr 1, tr  ? AB
Omr 2, vån ? 4
Omr 2, tr  ? DE
Omr 3, vån ? 1
Omr 3, tr  ? CEF
Omr 4, vån ? 3
Omr 4, tr  ? AD
Omr 5, vån ? 2
Omr 5, tr  ? E
Omr 6, vån ? 2
Omr 6, tr  ? D
Omr 7, vån ? 2
Omr 7, tr  ? BC
Omr 8, vån ? 3
Omr 8, tr  ? F

Svar: 24 
\end{verbatim}
\vfill
\columnbreak
\vspace{3cm}

%\begin{minipage}[c]{2\columnwidth+1\columnsep}
%    \centering\includegraphics[width=11cm]{orient.pdf}
%\end{minipage}
\vspace{0.5cm}
\;\\
\fe{Körningsexempel 1}\\
\begin{verbatim}
N ? 3
Omr 1, vån ? 2
Omr 1, tr  ? A
Omr 2, vån ? 100
Omr 2, tr  ? A
Omr 3, vån ? 3
Omr 3, tr  ? A

Svar: 195
\end{verbatim}
\columnbreak
%\vspace{8cm}
~\\
~\\
~\\
~\\
~\\
~\\
~\\
~\\
~\\

\ku{Figuren ovan:} Genomskärning av byggnaden i exempel 2. De stora siffrorna visar våningsnummer, de små numreringen av områdena. En möjlig snabbaste väg genom hela banan är:\\
%(1) AAD (2) EEE (3) CBAA (4) DEE (5) EEDD (6) DAAB (7) CFF (8)\\
AAD -- EEE -- CBAA -- DEE -- EEDD -- DAAB -- CFF \\
Notera att kontrollerna måste tas i ordning. %, så även om område 4 besökts redan på väg till område 2 så måste det besökas igen efter område 3.
\end{multicols*}


%På sin väg mellan två områden kan orienteraren naturligtvis behöva passera andra områden, men dessa måste ändå uppsökas på sin riktiga tid (eftersom orienteringskontroller alltid måste tas i en viss ordning).


\newpage
\specialsection*{Uppgift 4 -- Klippa protein}

Din kompis Björn har nyligen påbörjat ett stort kemiexperiment, med målet att skapa ett hemligt protein. 
Alla protein kan beskrivas som strängar av bokstäver från \texttt{a}-\texttt{z}, exempelvis \texttt{bcaa} och \texttt{hejsan}.
Han behöver skapa ett flertal så-kallade \textit{rena} proteiner. Dessa är protein som bara består av en enda bokstav.
För att göra detta har han mikroskopiska saxar som kallas CRISPR. Med hjälp av dessa kan han ta bort upp till $k$ bokstäver
som ligger bredvid varandra. Med andra ord kan du välja ett intervall med $k$ bokstäver att klippa bort.
Om ett protein delas i två delas när det klipps kommer dessa delar sedan att sättas ihop där de klipptes av. 

Eftersom han är väldigt lat vill han ha din hjälp att skriva ett program som beräknar minsta antalet gånger han behöver klippa ett visst protein för
att det ska bli rent. Notera att proteinet \textbf{inte} får bli en tom sträng, eftersom det isåfall inte längre kan användas för kemi.

Programmet ska först läsa in Björns protein, som består av som mest $32$ bokstäver. Sedan ska det läsa in talet $k$ ($k \ge 1$).

Programmet ska sedan skriva ut minsta antalet klippningar som krävs för att göra proteinet rent.

\vspace{0.5cm}
\fe{Poängsättning:}\\
De fem testfallen som vardera ger $1$ poäng har följande begränsningar utöver de som redan nämnts:
\begin{itemize}
    \item[(1)] $\;\;\;\;$ strängen består endast av bokstäverna \texttt{a} och \texttt{b}.    
    \item[(2)] $\;\;\;\;k=1$ 
    \item[(3-5)] $\;\;\;\; $Inga särskilda begränsningar.
\end{itemize}

\vspace{1cm}
\fe{Körningsexempel 1}\\
\begin{verbatim}
Björns protein ? exempelfall
k ? 4

Svar: 3
\end{verbatim}

\fe{Förklaring:}
Ett möjligt sätt är först klippa bort bokstäverna \texttt{pe}. Vi får då strängen \texttt{exemlfall}. \\
Sedan kan vi klippa bort \texttt{fa} och får \texttt{exemlll}. \\
Till sist klippa vi bort \texttt{exem} och får \texttt{lll}. 



\newpage
\specialsection*{Uppgift 5 -- Pickle clicker}

För ungefär ett halvår sedan släpptes succéspelet \emph{Pickle Clicker}. 
Efter att ha lagt otaliga timmar på spelet har Rasmus bestämt sig för att ta 
saker och ting till nästa nivå: han tänker speedrunna spelet. 


Pickle Clicker går ut på att man samlar ihop så mycket som möjligt av valutan pickles. 
Målet är att köpa Megapicklen så fort som möjligt, som kostar $t$ pickles. 
I spelet finns även $n$ olika typer av byggnader. De olika typerna är numrerade från 
$1$ till $n$, där den $i$:te typen producerar $p_i$ pickles per sekund och kostar 
$c_i$ pickles att bygga.

I början äger du en byggnad av typ $1$. Varje sekund sker följande: 
\begin{enumerate}
  \item Varje byggnad du äger producerar pickles. En byggnad av typ $i$ producerar $p_i$ pickles.
  \item Därefter har du alternativet att köpa en byggnad eller köpa Megapickeln. En byggnad av typ $i$ kostar 
  $c_i$ pickles att bygga, och du måste ha råd med den för att få köpa den. Du får maximalt köpa en byggnad 
  per sekund, men det är tillåtet att ha flera byggnader av samma typ.
\end{enumerate}


Hjälp Rasmus bestämma hur han ska spela spelet så att han kan köpa Megapicklen på så få sekunder som möjligt.
Se exemplet nedan för att få en tydligare bild av reglerna.

Programmet ska läsa in heltalen $n$ och $t$ ($1 \le n \leq 6, 1 \le t \le 10^5$): antalet byggnader och kostnaden av
Megapicklen.

Programmet ska skriva ut ett heltal: hur långa tid det tar att vinna pickle clicker om man spelar optimalt.

\fe{Poängsättning:}\\
De fem testfallen som vardera ger $1$ poäng har följande begränsningar utöver de som redan nämnts:
\begin{itemize}
    \item[(1)] $\;\;\;\;n=1$ 
    \item[(2)] $\;\;\;\;t \le 10$ 
    \item[(3)] $\;\;\;\;t \le 1000$
    \item[(4-5)] $\;\;\;\; $Inga särskilda begränsningar.
\end{itemize}


\fe{Körningsexempel 1}
\begin{verbatim}
Antal tärningar N ? 2
Antal sidor M ? 6 
Antal utfall K ? 2

Svar: 0.30555556
\end{verbatim}

\fe{Körningsexempel 2}
\begin{verbatim}
Antal tärningar N ? 3
Antal sidor M ? 7
Antal utfall K ? 4

Svar: 0.41399417
\end{verbatim}

\fe{Körningsexempel 3}
\begin{verbatim}
Antal tärningar N ? 20
Antal sidor M ? 2189
Antal utfall K ? 2734

Svar: 0.36902844
\end{verbatim}

\end{document}
